\subsection{The density clause and the size of the cubic families}
\label{sec:cubicdensityscope}

\begin{proposition}[P-F48: the cubic constructions are exponentially sparse in the stated input population]
\label{prop:cubicdensityscope}
For a positive integer $a$, put
$\ell(a)=1+\lfloor\log_2a\rfloor$. For a partition, sum this
quantity over its positive parts; zero parts are omitted.
Write
\[
 L(\lambda,\mu,\pi)=\ell(\lambda)+\ell(\mu)+\ell(\pi).
\]
This is exactly the input-size convention in GCT VI \cite{PGCTVI},
Section~1.1, printed p.~9. The narrow SH(b), Hypothesis~1.6.5(b),
printed p.~25, concerns triples of total length \emph{exactly}
$N$ with $a_{\lambda,\mu}^{\pi}\ne0$. The unified SH(b),
Hypothesis~3.3.2(b), printed p.~51, instead says length
\emph{at most} $N$ and uses the specified input $x$.
For its canonical ordinary-plethysm specification,
Section~3.4, printed p.~53, that length is
$n+L(\lambda,\mu,\pi)$: the group $GL_n$ contributes $n$.
Thus these are different input populations, not just
different cutoffs. Neither clause imposes a fixed inner
shape, fixed outer shape, fixed target shape, or a fixed rank.

Let $\mathcal D_N$ be the set of partition triples of length
exactly $N$ with nonzero ordinary plethysm coefficient, and
put $\mathcal D_{\le N}=\bigcup_{j\le N}\mathcal D_j$.
The counting fractions for narrow triples, using either cutoff,
and an exceptional set
$\mathcal E$ are
\[
 \frac{|\mathcal E\cap\mathcal D_N|}{|\mathcal D_N|},
 \qquad
 \frac{|\mathcal E\cap\mathcal D_{\le N}|}
      {|\mathcal D_{\le N}|},
 \tag{P-F48a}
\]
The first is the narrow clause's natural
conditional reading; the second is its at-most-$N$ variant.
The source states the density in words rather than displaying
a fraction. If instead that wording is read
as density among all partition triples, the estimates below
continue to hold, since that denominator is larger. No
distribution on partitions of a fixed weight or on a chosen
one-parameter family is substituted for the source's population.

For $n\ge4$, set
\[
 g=\gcd(n,3),\quad m=\binom{n+2}{3},\quad
 d=\frac{rn}{g},\quad k=\frac{3r}{g},
 \qquad r\ge1.
\]
The rescaled cubic triples
\[
                  ((d),(3),(k^n))
 \tag{P-F48b}
\]
form a set $\mathcal C$. Their stretching functions are
$f_n(rs)$, in the notation of P-F31--P-F32.
More generally, for arbitrary integers $q,b\ge0$ put
\[
 \begin{split}
  \mu_{n,q}&=(q+3,q^{\,n-1}),\\
  \lambda_{n,r,b}&=(d+b,b^{\,m-1}),\\
  e_{n,q}&=\frac{m(3+nq)}n
          =\frac{(n+1)(n+2)}2+qm,\\
  \pi_{n,r,q,b}&=((k+qd+be_{n,q})^n),
 \end{split}
 \tag{P-F48c}
\]
again omitting all zero parts. These natural inner and outer
determinant twists form a set $\mathcal C^{\det}$ containing
$\mathcal C$. Every one of them has the same stretching
function $f_n(rs)$ as its untwisted counterpart.

For every integer $N\ge12$,
\[
 \begin{split}
 |\mathcal D_N|&\ge 2^{N/2-4},\\
 |\{x\in\mathcal C:L(x)\le N\}|
       &\le N\,2^{N/5},\\
 |\{x\in\mathcal C^{\det}:L(x)\le N\}|
       &\le N(N+1)^3\,2^{N/4}.
 \end{split}
 \tag{P-F48d}
\]
Consequently the conditional density of \emph{any} subset
of $\mathcal C$ is at most
$16N\,2^{-3N/10}$, and that of any subset of
$\mathcal C^{\det}$ is at most
$16N(N+1)^3\,2^{-N/4}$, for both populations in
(P-F48a). In particular these densities are smaller than
$N^{-c}$ for every fixed $c>0$ and all sufficiently large
$N$.

The same conclusion holds separately in the canonical
ordinary-plethysm input population of the unified clause.
Let $\widehat{\mathcal D}_{\le N}$ count admissible
$(h,\lambda,\mu,\pi)$ with nonzero coefficient and
$h+L(\lambda,\mu,\pi)\le N$. Lift the constructed
families using all their admissible ambient ranks
$h\ge n$, including padding above the displayed rank $n$.
For $N\ge13$ their conditional densities in this population
are at most
\[
 32N^2\,2^{-3N/10},\qquad
 32N^2(N+1)^3\,2^{-N/4},
 \tag{P-F48h}
\]
respectively. This statement uses the canonical
ordinary-plethysm specification; it is not a claim
about counting arbitrarily many syntactic descriptions
of a composite subgroup embedding.

The unscaled family $r=1$, including all its twists in
(P-F48c), has zero initial multiplicity and is excluded
from $\mathcal D_N$ altogether. These facts do not prove
SH(b), do not assert that the rescaled functions have zero
strict saturation index, and do not preclude a different
family of positive density.
\end{proposition}

\begin{proof}
\emph{1. A nonzero population of exponential size in every
large length shell.}
For odd $N=2u+1$, choose any integer
$a\in[2^{u-1},2^u-1]$ and take
\[
                         (\lambda,\mu,\pi)=((1),(a),(a)).
\]
Its coefficient is one, since the outer Schur functor is
the identity. Its length is exactly $1+u+u=N$.
There are $2^{u-1}$ such triples.

For even $N=2u+4$, choose
\[
                    \left\lceil\frac{2^{u+1}}3\right\rceil
                              \le a\le 2^u-1
\]
and take $((3),(a),(3a))$. Here $\ell(a)=u$ and
$\ell(3a)=u+2$, so the total length is $2+u+(u+2)=N$.
The coefficient is one: the highest weight $(3a)$ occurs
once in $\operatorname{Sym}^3(\operatorname{Sym}^a V)$.
Indeed its weight space is spanned by the cube of the
highest weight vector of $\operatorname{Sym}^a V$.
This also gives a direct verification in any sufficient
ambient rank, or just in rank one.
For $u\ge4$, the number of choices is at least
$2^u/3-1\ge2^{u-2}$.
Both parity cases therefore give the first inequality in
(P-F48d) for $N\ge12$. In particular the denominator is
nonzero in every such shell.

\emph{2. The determinant twists preserve the actual ray.}
Let $V=\mathbb C^n$, $W=\operatorname{Sym}^3 V$, and
$W_q=W\otimes(\det V)^q$. Then
\[
              W_q=V_{\mu_{n,q}}(GL_n),\qquad \dim W_q=m.
\]
The determinant character of $W$ is
$(\det V)^{3m/n}$. To see this, each diagonal coordinate
occurs with the same total exponent in its monomial
basis, and the sum of those $n$ exponents is $3m$.
The integer $3m/n=(n+1)(n+2)/2$ proves both the formula
and the asserted integrality of $e_{n,q}$.
Consequently
\[
                    \det W_q=(\det V)^{e_{n,q}}.
\]
For every integer $s\ge0$, the usual determinant twist
identity for a highest weight gives
\[
 \begin{split}
 V_{s\lambda_{n,r,b}}(GL(W_q))
   &\simeq \operatorname{Sym}^{sd}(W_q)
                          \otimes(\det W_q)^{sb}\\
   &\simeq \operatorname{Sym}^{sd}(W)
                    \otimes(\det V)^{s(qd+be_{n,q})}.
 \end{split}
 \tag{P-F48e}
\]
The target $V_{s\pi_{n,r,q,b}}(GL_n)$ is the character
$(\det V)^{s(k+qd+be_{n,q})}$. Cancelling the common
determinant twist reduces its multiplicity to that of
$(\det V)^{sk}$ in $\operatorname{Sym}^{sd}(W)$.
The central degree condition is $3sd=nsk$, so that
this is precisely the $SL_n$-invariant multiplicity
$f_n(rs)$. This proves the ray identity for all $s$,
including the boundary cases $q=0$ and $b=0$.

\emph{3. The bit budgets are literal partition budgets.}
For (P-F48b), $d\ge r$ and $k\ge r$, whence
\[
             L((d),(3),(k^n))
                      \ge 2+(n+1)\ell(r).
 \tag{P-F48f}
\]
The target has $n$ positive parts, so $n\le N$ whenever
the total length is at most $N$. For each such $n$,
the number of possible positive $r$ is at most
$2^{N/(n+1)}\le2^{N/5}$.
Summing this bound over at most $N$ values of $n$
proves the second inequality in (P-F48d).

For the larger family define $\ell(0)=0$ only for the
two parameters $q,b$. This convention counts their
omitted zero parts; it does not add zeros to the source's
partition encoding. Put
$u=\ell(r)$, $v=\ell(q)$ and $w=\ell(b)$.
The first part of $\lambda_{n,r,b}$ is at least $r$;
its other $m-1$ positive parts, when $b>0$, equal $b$.
The $n$ parts of $\mu_{n,q}$ are each at least $q$
when $q>0$. The target has $n$ parts each at least $r$.
These are disjoint entries in the three encoded
partitions, and therefore
\[
              L(\lambda_{n,r,b},\mu_{n,q},\pi_{n,r,q,b})
                   \ge (n+1)u+nv+(m-1)w
                   \ge4(u+v+w).
 \tag{P-F48g}
\]
When $q=0$ the omitted contribution on the right is
zero, while the actual inner partition still has its
positive part $3$; when $b=0$ only the first outer
part remains. Thus (P-F48g) is valid also in those
cases without double counting any entry.

There are at most $(N+1)^3$ possible bit-length triples
$(u,v,w)$. For each one, the number of parameter triples
$(r,q,b)$ is at most $2^{u+v+w}\le2^{N/4}$, using the
single zero option when $v=0$ or $w=0$.
Again $n\le N$ because the target has $n$ positive
parts. Summing gives the third inequality in
(P-F48d). Distinct parameter choices that happen to
give the same input only reduce the count.

\emph{4. The density and nonzero-input conclusions.}
For the exact-length density, bound its numerator by
the corresponding at-most-$N$ count in (P-F48d) and
its denominator by the first line there.
For the at-most-$N$ narrow density use the same numerator
bound and $|\mathcal D_{\le N}|\ge|\mathcal D_N|$.
This gives exactly the two displayed exponential
estimates, even if every member of the constructed
family were an SH(b) exception. Exponential decay
times a fixed polynomial is smaller than every
inverse power of $N$.

For the canonical unified ordinary-plethysm inputs,
each underlying triple has at most $N$ admissible
ambient ranks of cost at most $N$.
The stable plethysm coefficient, including its stretched
versions, is unchanged by such padding: specialization
of the symmetric-function character to $h$ variables
preserves the coefficients of all Schur functions of
height at most $h$, including the target.
Multiplying the previous numerator bounds by $N$
therefore covers every padded realization of the family.
The one-row examples in the first part of the proof
are all admissible already at rank one.
Use those of narrow length $N-1$, then add the rank
cost one. They show
\[
        |\widehat{\mathcal D}_{\le N}|
                  \ge 2^{(N-1)/2-4}
                  \ge 2^{N/2-5}
                         \qquad(N\ge13).
\]
Combining this with the enlarged numerator bounds
proves (P-F48h). Any larger population with these
canonical inputs and the same encoding costs has a
larger denominator; no inference about a different
encoding's numerator is made.

Finally P-F32 applied at $s=1$ gives
$f_n(1)=0$ for every $n\ge4$: its exterior parameter
is $k=3/\gcd(n,3)\le3$, for which
$(k-1)(k-2)/2\le1<n$.
Equation (P-F48e) preserves this zero under all the
determinant twists.
If instead $f_n(r)>0$, powers of a nonzero invariant
show $f_n(rs)>0$ for every positive integer $s$.
This last statement concerns realized values only.
Definition~1.2.4 of GCT VI tests every nonzero
constituent polynomial at every positive integer,
including arguments outside that constituent's
residue class. Realized nonvanishing alone thus
does not establish the strict saturation assertion,
and none was used in the counting proof.
\end{proof}

\paragraph{Scope of the remaining density problem.}
The two source versions use different length cutoffs
and, in their canonical ordinary-plethysm specifications,
different rank costs. The proof explicitly treats
each, so their equivalence is not assumed.
Neither changing the conditioning to
rectangular targets nor counting only the parameter
$r$ would test the stated clause. Nor does padding an
ambient rank create new narrow inputs: the source
counts the three partitions and omits their zero parts.
It does create distinct canonical unified inputs,
and the extra factor $N$ in (P-F48h) explicitly covers
them.
The proved conclusion closes this particular
amplification attempt, including its natural determinant
twists. Establishing or refuting SH(b) for the entire
source population remains open here.
